\section{Introduction}

Reproducible builds provide an important foundation for trust in released software~\cite{lamb_reproducible_2022,fourne_its_2023,benedetti_empirical_2025}: given the same source and build instructions, an independent party should produce an identical artifact, detecting build-time tampering and confirming that a released artifact corresponds to the reviewed source~\cite{WhyReproducibleBuilds,BestpracticesSoftwaresupplychainOssscbestpracticesmd}. Open-source communities, standards, and government agencies accordingly recommend them as a supply-chain defense~\cite{moore__software_2024,CISANSAODNI2022,sigOpenSourceProject}---especially for ecosystems such as npm and PyPI, which face a growing wave of attacks on compromised build and release pipelines~\cite{snykNodeIpc2026,snykAntV2026,stepsecurityDurabletask2026,snykTanstack2026}.

However, reproducibility alone is insufficient for verifying the integrity of released artifacts~\cite{xiongBuildVerifiabilityJavabased2022,decarnedecarnavaletChallengesImplicationsVerifiable2014}.
An independent verifier must also know the source revision, build environment, and build instructions that produced the artifact in order to reproduce it.
Centralized-build ecosystems such as Debian~\cite{lamb_reproducible_2022} and Nix~\cite{malkaDoesFunctionalPackage2025} make this easier by preserving or publishing artifact-bound build specifications through ecosystem infrastructure.
In contrast, \textit{decentralized-build} ecosystems rely on third-party build infrastructure and diverse release practices, leading to fragmented or missing source and build metadata~\cite{goswami_npm_2020, gao_pyradar_2024}.
Prior work has largely avoided this problem by assuming the relevant metadata is available~\cite{benedetti_empirical_2025}, relying on human augmentation~\cite{google_oss_rebuild}, or treating artifact verification as out of scope~\cite{hassanshahiUnlockingReproducibilityAutomating2025}.
As a result, it remains unclear whether independent artifact verification is feasible in source-oriented package ecosystems.


This paper addresses this gap empirically. We focus on decentralized-build package ecosystems, such as npm and PyPI, because the growing prevalence of software supply-chain attacks highlights the urgency of reproducible builds and verifiability in these ecosystems~\cite{okafor_sok_2022,sonatype_state_2026}.
% This paper addresses this gap empirically.
% We focus on decentralized-build package ecosystems, such as npm and PyPI, because the growing wave of supply-chain attacks highlight the urgency of reproducible builds and verifiability in these ecosystems. 
We define two models for evaluating verifiability: an independent verifier model using only registry-exposed and registry-derived metadata and an artifact comparison model that allows graded levels of artifact equivalence.
We implement these models in \toolname, an Artifact Verification Pipeline that recovers rebuild metadata, reconstructs build environments, rebuilds artifacts, and compares rebuilt outputs against registry artifacts using tiered equivalence levels. 
We use \toolname to measure verification rates, failure causes, and the effects of provenance attestations, artifact characteristics, and development choices across the four ecosystems: Crates.io, npm, PyPI, and RubyGems.


Our results show that artifact verifiability is low and uneven: across ecosystems, roughly 53–87\% of completed rebuilds verify, and recovering the source commit is the dominant obstacle. Provenance attestations improve recovery, but their benefit comes mostly from the development practices of the projects that adopt them rather than the recorded metadata itself; most remaining mismatches trace to unpinned build-tool and dependency versions, CI/CD-generated files, and source modifications, and many are small or metadata-only---byte identity alone is too coarse for source-oriented artifacts. How a package is \emph{built}, its packaging determinism, predicts verifiability far better than what it \emph{is}, with native compilation a structural limit. Finally, against the state-of-practice OSS-Rebuild, \toolname verifies about twice as many artifacts and, unlike curated rebuild platforms, scales to the popular long tail.

In summary, this paper makes three contributions:
\begin{itemize}
\item We define new models for independent artifact verifiability and implement them in \toolname for recovering rebuild metadata, rebuilding and verifying published software artifacts.
\item We empirically study the verifiability of 34{,}005 packages across four source-oriented package ecosystems.
\item We identify the metadata gaps, release practices, and ecosystem-level changes that affect verifiability.
\end{itemize}

\textbf{Significance:}
This paper provides the first cross-ecosystem measurement of independent artifact verifiability. 
Verifiability in decentralized-build ecosystems remains limited by insufficient rebuild metadata. By identifying the metadata gaps and packaging practices that prevent verification, the study provides guidance for registries, standards, and tooling to support verifiability as a practical software supply-chain defense.




% \section{Introduction}

% Modern software systems are assembled from large collections of third-party components. 
% Although package managers simplify the integration of third-party code, dependency adoption and updates are not necessarily accompanied by complete source-code review.
% Prior studies have found that dependency updates are frequently only partially reviewed and that developers often do not manually review dependencies before adding or updating them~\cite{imtiaz_dependencies_reviewed_2023,hamer_trusting_code_2025}.

% These registries distribute the files that are ultimately installed into downstream applications: Python source distributions and wheels, npm tarballs, and Rust crate archives.
% Consequently, trust in open-source software depends not only on whether source code is publicly available, but also on whether the artifact distributed by a registry can be independently related back to that source.

% % Package managers are becoming more critical infrastructure in the modern software supply chain.
% % Developers rarely build software they use from scratch; instead, they depend on a large and growing 
% % ecosystem of third-party packages distributed through language-specific registries. In Python, the dominant 
% % distribution channel is Python Package Index (PyPI), where users consume source distributions and wheels rather than source repositories directly.
% % As a result, trust in open-source software increasingly depends not only on the availability of source code, but also on 
% % whether the released artifact on the registry can be independently related back to that source.

% % Reproducible builds (R-Bs) provide an important foundation for this problem~\cite{lamb_reproducible_2022,fourne_its_2023,benedetti_empirical_2025}.
% % In principle, if an independent party is given the correct source revision, the build instructions, and a sufficiently similar build environment, they should 
% % be able to rebuild the same artifact. This property helps detect build-time tampering and strengthens confidence that a released package corresponds to the 
% % source code that users and auditors inspect. However, reproducibility alone does not solve the full verification problem for package registries.
% % To independently confirm a released artifact, one must first identify the correct public source repository, recover the specific release commit, reconstruct 
% % enough of the publisher's build environment, and only then compare the rebuilt artifact against the published one.

% Reproducible builds (R-Bs) provide an important foundation for establishing this correspondence~\cite{lamb_reproducible_2022,fourne_its_2023,benedetti_empirical_2025}.
% Given the correct source revision, build instructions, and a sufficiently controlled build environment, an independent party should be able to produce an equivalent artifact.
% This property helps detect build-time tampering and strengthens confidence that a released package corresponds to the source code reviewed by users and auditors.
% Prior work has shown that reproducibility varies substantially across packaging ecosystems, but that many failures can be addressed through relatively small changes to packaging tools and infrastructure defaults~\cite{benedetti_empirical_2025}.
% However, reproducibility alone does not solve the end-to-end verification problem for public package registries.
% It assumes that the verifier already knows which repository, source revision, packaging workflow, and environment produced the published artifact.

% Recovering this information is itself a non-trivial problem.
% Prior studies have identified discrepancies between PyPI distributions and their associated repositories~\cite{vu_lastpymile_2021}, examined the reproducibility of npm packages against upstream source~\cite{goswami_npm_2020}, and shown that registry metadata may omit or misidentify source repositories~\cite{PypiNpmGHAccessibility,PyRadar}.
% These studies establish important parts of the problem, but they leave an unresolved last-mile question:
% Starting from an artifact distributed by a public registry, can an external party recover sufficient source and build information to independently confirm its correspondence to a public source revision?


% This question is broader than byte-for-byte reproducibility alone.
% Prior work distinguishes deterministic build processes, which seek to eliminate sources of nondeterminism, from explainable build processes, which identify and interpret residual differences that cannot or should not be eliminated~\cite{shi_experience_2022}.
% For package registries, a mismatch between a rebuilt artifact and its published counterpart may arise from benign archive metadata, timestamps, generated files, ecosystem-specific packaging transformations, or environment-dependent behavior.
% Alternatively, it may indicate that the verifier selected the wrong source revision, reconstructed the wrong workflow, or encountered an unexplained divergence between the public source and the artifact distributed by the registry.
% Independent artifact verification is therefore a multi-stage problem rather than a single hash comparison.

% This problem is especially timely because package ecosystems increasingly expose verification-relevant evidence, but do so in different forms.
% PyPI supports index-hosted attestations associated with individual release files through PEP~740~\cite{pep740,pypi-integrity-api,pypi-publish-attestation}.
% npm exposes provenance statements that can associate a published package with its build environment, source commit, and publishing workflow~\cite{npm-provenance}.
% Cargo embeds a VCS info \iffalse \texttt{.cargo\_vcs\_info.json} \fi file in packaged crate archives, recording the source checkout hash when available and the package's repository-relative path~\cite{cargo-package}.
% In addition, crates.io supports OIDC-based Trusted Publishing for CI/CD workflows~\cite{crates-trusted-publishing}.
% These mechanisms provide valuable evidence, but they differ in coverage, semantics, and completeness.
% In particular, evidence that identifies a publishing workflow does not necessarily provide the complete build recipe required to reconstruct an artifact or establish source--artifact equivalence.

% In this paper, we study the practical verifiability of artifacts distributed through PyPI, npm, and crates.io.
% We define \emph{artifact verifiability} as the ability of an external party to start from a released registry artifact, obtain or recover sufficient public evidence to identify its source revision and packaging process, rebuild the artifact in a reconstructed environment, and compare the rebuilt output against the published counterpart.
% We distinguish among multiple evidence paths: direct ecosystem-provided evidence, heuristic recovery, and their combination.
% We then analyze where verification succeeds and where it breaks down.

% To support this analysis, we develop \toolname, an automated Artifact Verification Pipeline.
% \toolname recovers source and build metadata, reconstructs ecosystem-specific packaging environments, rebuilds each artifact, and compares the rebuilt output against the registry artifact using explicitly defined equivalence levels.
% \toolname also distinguishes failures of source recovery, build reconstruction, internal build repeatability, and artifact equivalence.
% Using representative samples from PyPI, npm, and crates.io, we measure the extent to which publicly distributed artifacts can be independently traced to source revisions and rebuilt across three packaging ecosystems with substantially different metadata, tooling, and distribution models.
